\documentclass[12pt,a4paper]{article}

\usepackage{amsmath, amssymb}
\usepackage{geometry}
\usepackage{parskip}
\usepackage{times}

\geometry{top=2cm, bottom=2cm, left=2.5cm, right=2.5cm}
\linespread{1.1}
\setlength{\parindent}{0pt}
\setlength{\parskip}{3.5pt}

\begin{document}

% ================================================================
\subsection*{Self-Exciting Hawkes Jump Models}
% ================================================================

% ----------------------------------------------------------------
\subsubsection*{Introduction to the Paper}
% ----------------------------------------------------------------

A key challenge in modelling oil prices is that large price
movements do not arrive independently. In practice, a sudden
crash tends to trigger further extreme moves shortly afterwards
--- a feature known as \textit{jump clustering}. Standard
jump-diffusion models handle jumps through a Poisson process with
a \textit{fixed} arrival rate, which cannot reproduce this
behaviour. Gonzato and Sgarra (2020) address this by proposing a
\textit{self-exciting} jump-diffusion model where the jump
intensity is stochastic and path-dependent: it rises after each
jump and gradually decays back to a long-run background level.
Combined with stochastic volatility and a stochastic convenience
yield, the model is far more realistic than earlier factor models
and is shown to improve both in-sample fit and out-of-sample
futures forecasting for WTI crude oil data from 2008 to 2018.

% ----------------------------------------------------------------
\subsubsection*{What the Paper Introduced}
% ----------------------------------------------------------------

The key building block is the \textit{Hawkes process}, whose
conditional intensity $\lambda_t$ satisfies:
\begin{equation}
    d\lambda_t = \beta(\lambda_{\infty} - \lambda_t)\,dt
    + \alpha\,dN_t,
\end{equation}
where $\lambda_\infty>0$ is the background intensity, $\beta>0$
controls the speed of decay after a jump, and $\alpha>0$ is the
magnitude of the self-exciting spike caused by each new event
$dN_t$. Every jump immediately raises $\lambda_t$ by $\alpha$,
making the next jump more likely; the effect then fades
exponentially, so calm periods are also possible. This reproduces
\textit{jump clustering} directly through the model structure.
The log-price $x_t = \ln(S_t/S_0)$ follows:
\begin{equation}
    dx_t = \!\left(\mu - \tfrac{1}{2}V_t
    - \lambda_t\mathbb{E}[e^{J_x}-1]
    - \delta_t\right)dt
    + \sqrt{V_t}\,dW_t + dJ_{x,t},
\end{equation}
where $V_t$ is a CIR-type stochastic variance that co-jumps with
the price, and $\delta_t$ is an Ornstein--Uhlenbeck convenience
yield. The affine structure of the system allows futures prices
to be computed in closed form, greatly simplifying estimation.

% ----------------------------------------------------------------
\subsubsection*{Why Hawkes Processes Matter}
% ----------------------------------------------------------------

In a standard Poisson model, the jump probability on any given
day is independent of the past, which is clearly unrealistic.
After a major oil market shock --- a geopolitical event, a supply
disruption --- uncertainty stays elevated and further large moves
remain more likely for days or weeks. Hawkes processes model this
\textit{financial contagion} directly: each jump raises the
intensity, so consecutive extreme events are more probable.
This makes them far better suited than constant-intensity models
for capturing the turbulent, clustered dynamics that dominate
oil markets during crises, and for producing realistic estimates
of tail risk and derivative prices.

% ----------------------------------------------------------------
\subsubsection*{Bayesian Estimation and Sequential Monte Carlo}
% ----------------------------------------------------------------

Estimation is complicated by the fact that $V_t$, $\delta_t$,
and $\lambda_t$ are all \textit{latent states} that cannot be
observed directly. The authors use the \textit{two-stage
Sequential Monte Carlo (SMC) sampler}: a swarm of weighted random
samples (\textit{particles}) approximates the posterior of the
parameters and latent states, being gradually tempered from the
prior to the posterior via adaptive importance sampling and
resampling. A first stage with $M_1=30$ state particles explores
the parameter space cheaply; a second stage with $M_2=500$
particles corrects the approximation. The method also delivers
the marginal likelihood, enabling formal Bayes-factor model
comparison --- something that would be very difficult with
classical maximum-likelihood approaches given the non-linear,
non-Gaussian state-space structure.

% ----------------------------------------------------------------
\subsubsection*{Empirical Results}
% ----------------------------------------------------------------

Applied to 2{,}767 daily WTI observations (spot plus eight
futures contracts, 2008--2018), the results strongly support
self-excitation. Log-returns exhibit kurtosis of $7.31$ and fail
the Jarque--Bera normality test. \textbf{Model~I} (Hawkes
dynamics) is compared to \textbf{Model~II} (constant intensity,
$\alpha=\beta=0$): the log Bayes factor equals $9.13$,
constituting \textit{decisive} evidence for Model~I. The
self-excitation parameter is estimated at $\hat{\alpha}=135.19$
(SD~$5.02$), confirming that jump clustering is a genuine feature
of the oil market. Out of sample, Model~I also achieves lower
Mean Absolute Error in futures pricing across all maturities.

% ----------------------------------------------------------------
\subsubsection*{Strengths of the Paper}
% ----------------------------------------------------------------

The paper models jump dynamics in a flexible and realistic way,
with the Hawkes intensity adapting automatically to market
conditions. Stochastic volatility with co-jumps captures the
leverage effect, and the affine structure keeps futures pricing
analytically tractable. The model also has clear practical value:
two optimal time-varying hedge ratios are derived, and the
skewness-adjusted variance--skewness strategy outperforms pure
minimum-variance hedging both in and out of sample.

% ----------------------------------------------------------------
\subsubsection*{Limitations of the Paper}
% ----------------------------------------------------------------

The main drawbacks are computational. Running the two-stage SMC
sampler is resource-intensive, and the large number of parameters
increases calibration difficulty and the risk of overfitting.
Results may also be somewhat sensitive to the choice of prior
distributions, even though the authors report low prior
sensitivity for most parameters. Overall, the Hawkes framework
is considerably harder to implement than a standard
constant-intensity model, which may limit its practical adoption.

% ----------------------------------------------------------------
\subsubsection*{Connection to Our Project}
% ----------------------------------------------------------------

This paper is directly relevant because technology-heavy markets
like QQQ also exhibit clustered extreme movements: crashes tend
to arrive in waves rather than as isolated events. The Merton
model assumes constant Poisson intensity and therefore
underestimates the probability of consecutive large moves, while
the Kou model improves jump asymmetry but keeps the intensity
fixed. Hawkes-type models represent a natural frontier extension,
allowing the intensity to respond dynamically to market shocks.
Although implementing a full SMC-based self-exciting model is
beyond the scope of this project, understanding how Hawkes
processes extend traditional jump-diffusion frameworks provides
important context for interpreting the results obtained with the
Merton and Kou models and motivates the modelling choices made
throughout the report.

\end{document}